\section{Preliminary Results}
All results in this section are preliminary unless explicitly described as a completed correctness gate.

\subsection{Correctness and time to depth}
The Phase One release-scale rule, movement, wall-legality, path, and evaluation test completed \ParityPositions deterministic random comparisons with zero mismatches. Ten curated fixtures also matched through exhaustive depth three. These tests establish parity over the sampled state space, not a proof covering every reachable state.

At five seconds from the opening, the selective lane advanced from \PhaseOneOldFiveSecondSelectiveDepth{} to \PhaseOneNewFiveSecondSelectiveDepth{} completed plies, while the exhaustive lane advanced from \PhaseOneOldFiveSecondVerifiedDepth{} to \PhaseOneNewFiveSecondVerifiedDepth{}. At fifteen seconds, selective depth advanced from \PhaseOneOldFifteenSecondSelectiveDepth{} to \PhaseOneNewFifteenSecondSelectiveDepth{}. \Cref{fig:time-depth} plots these opening measurements.

\begin{figure}[t]
\centering
\begin{tikzpicture}
\begin{axis}[
  ybar,bar width=5pt,width=\columnwidth,height=4.2cm,
  ylabel={completed ply},symbolic x coords={5s-selective,5s-verified,15s-selective,15s-verified},
  xtick=data,x tick label style={rotate=28,anchor=east,font=\scriptsize},
  ymin=0,ymax=12,legend style={font=\scriptsize,at={(0.5,1.02)},anchor=south,legend columns=2},
  grid=major
]
\addplot table[x=case,y=before,col sep=comma]{generated/time-to-depth.csv};
\addplot table[x=case,y=after,col sep=comma]{generated/time-to-depth.csv};
\legend{before,Phase One}
\end{axis}
\end{tikzpicture}

\caption{Opening completed depth before and after the Phase One exact hot-path changes. Selective and exhaustive depths measure different coverage.}
\label{fig:time-depth}
\end{figure}

\subsection{Parallelism and reuse}
On the recorded i5-12600KF host, a two-second root-split screen measured the aggregate results in \cref{fig:workers}. These figures do not imply superlinear hardware efficiency because split searches visit different nodes. The dynamic exact root queue reduced a depth-five opening search from \DynamicRootSingleMs{} to \DynamicRootEightMs{} ms with eight workers, a \DynamicRootSpeedup$\times$ speedup and \DynamicRootEfficiency\% efficiency. It remains below the 60\% promotion target.

\begin{figure}[t]
\centering
\begin{tikzpicture}
\begin{axis}[
  width=\columnwidth,height=4cm,xlabel={workers},ylabel={aggregate NPS},
  xtick={1,4,12},scaled y ticks=base 10:-6,y tick scale label code/.code={$\times10^6$},
  ymin=0,grid=major
]
\addplot+[mark=*,thick] table[x=workers,y=nps,col sep=comma]{generated/worker-scaling.csv};
\end{axis}
\end{tikzpicture}

\caption{Aggregate root-worker NPS. Wall-clock completed depth, rather than cross-layout NPS, remains authoritative.}
\label{fig:workers}
\end{figure}

Persistent analysis improved the expected-move opening from \ColdExpectedMoveDepth{} to \WarmExpectedMoveDepth{} completed plies at one second and reused \ExpectedMoveReusedNodes{} nodes. Repeating the same opening reported \SamePositionReusedNodes{} reused nodes. These are cache-reuse measurements, not additional newly searched nodes.

\subsection{Self-play pilot}
The resource-capped pilot completed \PilotGames games and \PilotPlies plies, averaging \PilotAveragePlies{} plies per game. The hybrid won \PilotHybridWins, the exhaustive reference won \PilotExhaustiveWins, and \PilotUnresolved games reached the limit. \Cref{fig:pilot-outcomes} separates wins by engine and color.

\begin{figure}[t]
\centering
\begin{tikzpicture}
\begin{axis}[
  ybar,bar width=10pt,width=\columnwidth,height=4cm,
  ylabel={wins},symbolic x coords={periwinkle,blossom},xtick=data,
  legend style={font=\scriptsize,at={(0.5,1.02)},anchor=south,legend columns=2},
  ymin=0,grid=major
]
\addplot table[x=color,y=hybrid,col sep=comma]{generated/pilot-wins.csv};
\addplot table[x=color,y=exhaustive,col sep=comma]{generated/pilot-wins.csv};
\legend{hybrid,exhaustive}
\end{axis}
\end{tikzpicture}

\caption{Preliminary pilot wins by engine and color over both time controls. Unresolved games are omitted from the bars.}
\label{fig:pilot-outcomes}
\end{figure}

Neither 250 ms nor one-second SPRT crossed a configured boundary. Moreover, this older pilot allowed the color-swapped return game to inherit per-engine caches across a game boundary. The data remain useful as a baseline but cannot promote or reject the hybrid. The updated protocol clears each engine between games while preserving reuse between moves within a game.

\subsection{Negative and held experiments}
The experiment registry in \cref{tab:experiments} records failures rather than selecting only favorable measurements. Eager topology caching, partial selection, several untuned histories, canonical per-node symmetry, and the SIMD build reduced NPS or completed depth. ProbCut, multi-cut, reverse futility, adaptive reductions, strategic evaluation, quiescence, and all-shortest-route candidates remain held because short screens changed preferred moves or lacked strength evidence.

\subsection{Phase Three engineering tranche}
The first Phase Three tranche retained two exact hot-path changes and added same-root iteration continuation. The latter reports the restored completed depth explicitly and is bypassed by strict fixed-depth and fixed-node searches. It passed clear-state and strict-search checks, but its short fixed-time depth gain varied across repeated runs and is therefore not yet reported as a strength or performance improvement.

\Cref{tab:phase-three-pruning} records the first fixed-node selective screens. ProbCut scored \PhaseThreeProbCutWins--\PhaseThreeProbCutBaselineWins and multi-cut scored \PhaseThreeMultiCutWins--\PhaseThreeMultiCutBaselineWins. Both reached lower average depth than their baselines despite modestly higher reported NPS. They were rejected because faster node counting did not translate into deeper or stronger search. Adaptive LMR subsequently led \PhaseThreeAdaptiveLmrWins--\PhaseThreeAdaptiveLmrBaselineWins{} over \PhaseThreeAdaptiveLmrGames{} games and increased average depth from \PhaseThreeAdaptiveLmrBaselineDepth{} to \PhaseThreeAdaptiveLmrDepth{}. At that tranche, this was promising preliminary evidence below the planned game volume; the follow-up below adds fixed-time and accuracy audits.

% Generated file.
\begin{table*}[t]
\centering
\caption{Preliminary Phase Three fixed-node pruning screens. Each move received 50,000 nodes.}
\label{tab:phase-three-pruning}
\scriptsize
\begin{tabular}{lrrrrrrl}
\toprule
Candidate & Games & Cand. & Base & Unres. & Cand. depth & Base depth & Status \\
\midrule
ProbCut & 100 & 49 & 51 & 0 & 7.66 & 7.86 & rejected \\
multi-cut & 100 & 50 & 50 & 0 & 7.96 & 8.59 & rejected \\
reverse futility & 100 & 50 & 50 & 0 & 8.74 & 8.35 & held \\
adaptive LMR & 500 & 267 & 233 & 0 & 8.66 & 8.23 & held \\
\bottomrule
\end{tabular}
\end{table*}

\subsection{Selective-search validation}
A larger follow-up separated playing strength from fixed-depth accuracy.
Aggressive adaptive LMR scored \AggressiveLmrFixedNodeScore{} across
\AggressiveLmrFixedNodeGames{} fixed-node games, but increased summed
depth-five root regret from \SelectiveBaselineRegret{} to
\AggressiveLmrRegret{} over 50 random positions. It therefore remains held
despite its positive match score.

The guarded variant reduced summed regret to \GuardedLmrRegret{}, scored
\GuardedLmrFixedNodeScore{} over \GuardedLmrFixedNodeGames{} fixed-node games,
and scored \GuardedLmrTimedScore{} over \GuardedLmrTimedGames{} games at
250 ms per move. Those timed games averaged \GuardedLmrTimedAveragePlies{}
plies. Neither sequential test crossed a promotion boundary, and representative
timed benchmarks showed no completed-depth gain. A subsequent one-second
screen reversed the short-budget result: guarded LMR scored
\GuardedLmrOneSecondScore{} over \GuardedLmrOneSecondGames{} games, averaging
\GuardedLmrOneSecondAveragePlies{} plies at an aggregate
\GuardedLmrOneSecondNps{} NPS. Guarded LMR was rejected at this non-regression
gate, and the planned five-second screen was not run. The production mask
therefore remains zero.

% Generated file.
\begin{table*}[t]
\centering
\caption{Selective-search validation. Regret is summed over the same 50-position depth-five audit; lower is better. Both SPRTs remained undecided.}
\label{tab:selective-validation}
\scriptsize
\begin{tabular}{lrrr rrl}
\toprule
Variant & Fixed-node games & Fixed-node score & 250 ms games & 250 ms score & Regret & Status \\
\midrule
adaptive LMR & 2000 & 1018--981--1 & 100 & 64--35--1 & 1300 & held \\
guarded LMR & 500 & 264--236--0 & 100 & 62--38--0 & 1059 & held \\
\bottomrule
\end{tabular}
\end{table*}


\subsection{Topology and root scheduling}
The pawn-triggered topology design preserved exact fixture parity but reduced
one-second median NPS from \TopologyVFourBaselineNps{} to
\TopologyVFourCandidateNps{} and average completed depth from
\TopologyVFourBaselineDepth{} to \TopologyVFourCandidateDepth{}. It was
rejected because constructing two reverse fields still cost more than the
saved pawn-path searches.

The sixth dynamic-root prototype uses persistent warm workers, PV-first work,
score-volatility gating, and an exact serial fallback. At depth six, the stable
opening improved by \DynamicRootOpeningSpeedup$\times$ with
\DynamicRootWorkers{} workers, but parallel efficiency was only
\DynamicRootOpeningEfficiency\%. Unstable fixtures used the serial fallback;
an ungated run had previously exposed an equal-score tie divergence and severe
over-search. The guarded scheduler remains held below the 60\% gate.

The shared-memory Wasm artifact compiled and passed parity, but is not enabled.
Independent Emscripten instances cannot safely share their entire linear memory;
a native threaded runtime and explicitly shared caches remain future work.

% Generated file.
\begin{table*}[t]
\centering
\caption{Preliminary depth-six guarded dynamic-root results. Unstable positions use the exact serial fallback.}
\label{tab:phase-five-root}
\scriptsize
\begin{tabular}{llrrrr}
\toprule
Position & Mode & Serial ms & Scheduled ms & Speedup & Move parity \\
\midrule
opening & dynamic & 1803 & 977 & 1.85 & yes \\
channelled routes & serial fallback & 2414 & 2414 & 1.00 & yes \\
transposition rich & serial fallback & 2125 & 2125 & 1.00 & yes \\
\bottomrule
\end{tabular}
\end{table*}

\subsection{Exact TT quality and dynamic-root v7}
Phase 6 corrected clustered TT admission so a matching packed position is
found before an earlier empty slot can create a duplicate. Exact bounds receive
a small replacement preference. Timed search can use a shallow entry's move for
ordering without using its score; strict fixed-depth mode remains isolated from
this cross-search ordering signal. The standard parity screen again reported
zero mismatches.

The one-second ten-position screen preserved the aggregate exhaustive depth and
increased aggregate selective depth from \PhaseSixBaselineSelectiveDepth{} to
\PhaseSixCandidateSelectiveDepth{}. Mean NPS moved slightly downward in that
single run, while the 250 ms run moved slightly upward. The change is retained
for the exact cluster-capacity correction, but these wall-clock results remain
preliminary rather than evidence of a stable speed gain.

Dynamic-root v7 removes the speculative prescreen and first establishes one
exact incumbent. Alternatives are distributed with zero windows, while every
tie or fail-high receives a full-window re-search. The recorded eight-worker
opening improved by \PhaseSevenOpeningSpeedup$\times$ and reduced repeated root
searches from \PhaseSevenPriorRepeatedMoves{} to
\PhaseSevenRepeatedMoves{}. Its \PhaseSevenOpeningEfficiency\% efficiency
still failed the 60\% gate, and a two-worker run was slower than serial search.
Cheap or unstable iterations now use exact serial fallback. The scheduler
therefore remains held and is not connected to the browser or extension.

% Generated file.
\begin{table*}[t]
\centering
\caption{Preliminary depth-six seed-first dynamic-root results on the opening.}
\label{tab:phase-seven-root}
\scriptsize
\begin{tabular}{rrrrrrr}
\toprule
Workers & Serial ms & Scheduled ms & Speedup & Efficiency (\%) & Re-searches & Parity \\
\midrule
2 & 1695 & 1824 & 0.93 & 46.5 & 2 & yes \\
4 & 1663 & 1171 & 1.42 & 35.5 & 3 & yes \\
8 & 1724 & 838 & 2.06 & 25.7 & 4 & yes \\
\bottomrule
\end{tabular}
\end{table*}


\subsection{Shared-memory search and low-wall proofs}
Phase 8 replaced the earlier collection of isolated Wasm instances with an
experimental pthread build whose root searches use one shared linear memory
and one mutex-protected transposition table. Across six depth-four fixtures,
best move, score, and completed depth matched serial exhaustive search. Two
threads produced a median \PhaseEightTwoSpeedup$\times$ speedup at
\PhaseEightTwoEfficiency\% efficiency. Eight threads reached
\PhaseEightEightSpeedup$\times$ but only \PhaseEightEightEfficiency\%
efficiency because duplicated search grew faster than useful parallel work.
The implementation remains held pending extension-context recovery tests and
does not replace the isolated production pool.

Phase 9 added bounded exact AND/OR proof search whose certificates are replayed
independently against legal move generation. In one-total-wall near-goal tests,
it certified a win after a \PhaseNineWinDepth-ply proof while visiting
\PhaseNineWinStates{} states and a loss after a \PhaseNineLossDepth-ply proof
while visiting \PhaseNineLossStates{} states. These proof lengths are replayed
certificate depths, not optimal distances. General one- and two-wall
tablebases remain unfinished, and cycles or budget exhaustion return unknown.
